% Obsahový zdroj: import_hodiny/32_Permutace_bez_opakovani.tex a odpovídající PDF (obě strany).
% P1–P5: zdrojové Příklady 1–5; C1–C5: zdrojové Procvičování 1–5.
% Matematický obsah zachován; výklad, odvození a řešení pedagogicky rozvedeny.
% Upřesnění a opravy: viz KMB_29-41_Kontrola.md; původní import zůstává beze změny.
% Rozpis: 4_rocnik/tematicke_plany/Rozpis_vyuky_MAT_4_rocnik.tex, hodina 32.
% Plán: 4_rocnik/tematicke_plany/SPS.pdf, oddíly 2.1–2.3.
% Sazba: sablony/pripravy/TEMPLATE_Priprava.tex a společný priprava.sty.
% Vlastní doplnění: didaktický výklad, zdůvodnění, metodika; žádné nové číslované úlohy.
% Učebnice/09_Kombinatorika nemá obsahovou kapitolu; místní kombinatorické listy ani Realisticky.cz nejsou dostupné.
\documentclass[11pt,a4paper]{article}
% Spolecny styl pro ucitelske pripravy MAT (Work/MAT)
% Verze 1.1, 2026-10-23
\NeedsTeXFormat{LaTeX2e}
\ProvidesPackage{priprava}[2026/10/23 v1.1 Sablona ucitelskych priprav MAT]

% Zakladni typografie a cestina
\RequirePackage[T1]{fontenc}
\RequirePackage[utf8]{inputenc}
\RequirePackage[czech]{babel}
\RequirePackage{lmodern}
\RequirePackage{microtype}
\RequirePackage{amsmath,amssymb}
\RequirePackage{xcolor}
\RequirePackage{enumitem}
\RequirePackage[most]{tcolorbox}
\RequirePackage{tabularx,booktabs,array}
\RequirePackage{fancyhdr}
\RequirePackage{needspace}
\RequirePackage{geometry}
\RequirePackage[hidelinks]{hyperref}

\geometry{a4paper,top=17mm,bottom=18mm,left=23mm,right=23mm}
\definecolor{notegray}{gray}{0.39}
\setlength{\parindent}{0pt}
\setlength{\parskip}{2.5pt plus 1pt minus 0.5pt}
\setlist[itemize]{topsep=2pt,itemsep=1pt,parsep=0pt,leftmargin=1.45em}
\setlist[enumerate]{topsep=2pt,itemsep=2pt,parsep=0pt,leftmargin=1.65em}
\setcounter{secnumdepth}{0}
\renewcommand{\arraystretch}{1.12}

% Zahlavi: vlevo rocnik | globalni tema, vpravo cislo: nazev hodiny
% Zapati zustava zcela prazdne.
\pagestyle{fancy}
\fancyhf{}
\renewcommand{\headrulewidth}{0.35pt}
\renewcommand{\footrulewidth}{0pt}
\setlength{\headheight}{14pt}
\newcommand{\nastavPripravu}[4]{%
  \fancyhead[L]{\small #1\enspace\textbar\enspace #2}%
  \fancyhead[R]{\small #3: #4}%
}

% Minimalisticke cernobile boxy. Kazda definice zvlast.
\tcbset{commonlesson/.style={enhanced,breakable=false,colback=white,colframe=black,
  colbacktitle=white,coltitle=black,boxrule=0.45pt,
  arc=0.8mm,outer arc=0.8mm,left=3mm,right=3mm,top=1.6mm,bottom=1.6mm,
  toptitle=1mm,bottomtitle=0.7mm,fonttitle=\bfseries,
  attach boxed title to top left={xshift=2.5mm,yshift=-1.15mm},
  boxed title style={colback=white,colframe=white,boxrule=0pt,left=1mm,right=1mm}}}
\newtcolorbox{zadaniBox}[1][]{commonlesson,title={Zadání},#1}
\newtcolorbox{prikladBox}[1][]{commonlesson,title={Příklad},#1}
\newtcolorbox{definiceBox}[1][]{commonlesson,title={Definice},#1}

% Spolecny box pro zadani a jeho reseni, bez mezery mezi castmi.
% Pouziti:
%   \begin{prikladReseniBox}
%     Zadani ulohy.
%     \tcblower
%     \begin{align*} ... \end{align*}
%   \end{prikladReseniBox}
% Spodni cast ma jemny sedy podklad a delici linku; lze v ni sazet
% vice radku matematiky. Obe casti maji jen jeden spolecny ramecek.
\tcbset{solutionlesson/.style={
  commonlesson,
  enhanced jigsaw,
  breakable,
  colbacklower=black!3,
  segmentation style={solid,draw=black!25,line width=0.35pt},
  fontlower=\small\color{notegray},
  before lower={\noindent\textbf{Řešení:}\enspace},
  leftlower=3mm,rightlower=3mm
}}
\newtcolorbox{prikladReseniBox}[1][]{solutionlesson,title={Příklad},#1}
\newtcolorbox{zadaniReseniBox}[1][]{solutionlesson,title={Zadání},#1}

% Pomucky pro autora pripravy: sedive poznamky; \reseni pro starsi pripravy
\newcommand{\poznamka}[1]{\par\begingroup\small\color{notegray}\itshape
  \textnormal{\textbf{Poznámka pro učitele.}}\enspace #1\par\endgroup}
\newcommand{\reseni}[1]{\par\begingroup\small\color{notegray}
  \textbf{Očekávané řešení:}\enspace #1\par\endgroup}
\newcommand{\blok}[1]{\needspace{5\baselineskip}\subsection*{#1}}
\newcommand{\vyrok}[1]{\uv{#1}}

\nastavPripravu{4. ročník}{Kombinatorika}{32}{Permutace bez opakování}
\begin{document}


Vrať se k řazení kartiček. Tentokrát chceme počet určit i pro větší množinu, aniž bychom vypisovali všechna pořadí. Po uspořádání všech prvků přejdeme k obsazení jen několika rozlišených míst.
\poznamka{Předpoklady: pravidlo součinu, faktoriál a jeho krácení. Cíl: žák odvodí $n!$ a $n!/(n-k)!$, rozliší všechny prvky od části a zdůvodní blokovou i doplňkovou metodu. Podle importu a zadání patří úvod do variací bez opakování do této hodiny, přestože rozpis výslovně jmenuje pouze permutace.}

\blok{Výklad a odvození}
\begin{definiceBox}
\textbf{Permutace bez opakování} je každé uspořádání všech $n$ navzájem různých prvků. Jejich počet značíme $P(n)$, kde $n\in\mathbb N_0$.
\end{definiceBox}

Pro $n\geq1$ obsazujeme první místo $n$ způsoby, druhé $(n-1)$ způsoby až poslední jedním způsobem. Proto
\[
 P(n)=n(n-1)\cdots1=n!.
\]
Pro $n=0$ existuje jediné prázdné uspořádání, takže $P(0)=0!=1$. Záleží na pořadí, žádný prvek se neopakuje a použijeme všechny prvky.

Chceme-li, aby dva určité prvky sousedily, spojíme je do bloku. Každé přípustné pořadí odpovídá právě jednomu pořadí bloků a jedné vnitřní orientaci dvojice. U dvou různých osob jsou uvnitř bloku dvě pořadí; tento faktor nesmíme vynechat. Zákaz sousedství lze řešit odečtením všech pořadí se sousedstvím od všech pořadí.\begin{definiceBox}
\textbf{Variace bez opakování} délky $k$ z $n$ navzájem různých prvků je uspořádaný výběr $k$ různých prvků. Pro celá $n\geq0$, $0\leq k\leq n$ značíme jejich počet $V(n,k)$.
\end{definiceBox}

Nyní obsadíme jen $k$ rozlišených míst. Pro $k\geq1$ má první místo $n$ možností a poslední $(n-k+1)$, tedy
\[
 V(n,k)=n(n-1)\cdots(n-k+1)=\frac{n!}{(n-k)!}.
\]
Pro $k=0$ je výběr prázdný a $V(n,0)=1$. Pro $k=n$ dostáváme $V(n,n)=P(n)$. Při $k>n$ nemůžeme bez opakování žádný takový výběr sestavit; faktoriálový vzorec v tomto oboru nepoužíváme. V této sadě dodržujeme zápis $V(n,k)$ se základem $n$ na prvním místě.


\poznamka{Průběh: 0–9 min odvození a P1; 9–21 min P2–P3; 21–31 min variace a P4; 31–42 min C1–C4; 42–45 min shrnutí. P5 a C5 tvoří rezervu či domácí procvičení. Hodina je obsahově hutná: blokové úlohy nesmějí vytlačit pochopení variací. Při nedostatku času společně dokonči P2, P3 vysvětli krátce doplňkem.}

\blok{Společně řešené příklady}

% Zdroj P1: import_hodiny/32_Permutace_bez_opakovani.tex, Příklad 1.
\begin{prikladReseniBox}[title={Příklad P1}]
Kolika způsoby se do řady postaví šest různých studentů?
\tcblower
Jde o permutaci všech šesti osob: \(P(6)=6!=720\).
\end{prikladReseniBox}


% Zdroj P2: import_hodiny/32_Permutace_bez_opakovani.tex, Příklad 2.
\begin{prikladReseniBox}[title={Příklad P2}]
V šestičlenné řadě mají Adam a Bára stát vedle sebe. Kolik je pořadí?
\tcblower
Adam a Bára tvoří blok, se čtyřmi dalšími studenty tedy řadíme pět různých objektů. To lze $5!$ způsoby. V bloku jsou pořadí $AB$ a $BA$, proto $5!\cdot2=240$. Spojení i rozpojení bloku je jednoznačné; žádné přípustné pořadí nechybí ani není započítáno dvakrát.
\end{prikladReseniBox}


% Zdroj P3: import_hodiny/32_Permutace_bez_opakovani.tex, Příklad 3.
\begin{prikladReseniBox}[title={Příklad P3}]
Šest různých studentů včetně Adama a Báry postavíme do řady. Kolik pořadí splňuje, že Adam a Bára nestojí vedle sebe?
\tcblower
Všech $6!=720$ pořadí rozdělíme na dvě nepřekrývající se skupiny: Adam a Bára sousedí, nebo nesousedí. První skupina má podle blokové metody $2\cdot5!=240$ pořadí. Ve druhé tedy zbývá $720-240=480$ pořadí. Neodečítáme jen orientaci $AB$, ale i $BA$.
\end{prikladReseniBox}


% Zdroj P4: import_hodiny/32_Permutace_bez_opakovani.tex, Příklad 4.
\begin{prikladReseniBox}[title={Příklad P4}]
Osm různých týmů soutěží o první tři místa. Kolik je výsledků, nejsou-li dělená umístění a každý tým obsadí nejvýše jedno místo?
\tcblower
Rozlišujeme vítěze, druhý a třetí tým. Každý tým může být vybrán jen jednou: $8$ možností pro zlato, $7$ pro stříbro a $6$ pro bronz. Tedy
\[
 V(8,3)=8\cdot7\cdot6=336=\frac{8!}{5!}.
\]
Ostatních pět týmů se do výsledku neřadí. Záměna zlatého a stříbrného týmu mění výsledek.
\end{prikladReseniBox}


% Zdroj P5: import_hodiny/32_Permutace_bez_opakovani.tex, Příklad 5.
\begin{prikladReseniBox}[title={Příklad P5}]
Kolik sedmiciferných čísel zakončených sudou cifrou vytvoříme z číslic $1,\ldots,7$ bez opakování?
\tcblower
Na posledním místě jsou \(3\) možnosti: \(2,4,6\). Zbývajících šest číslic lze uspořádat \(6!\) způsoby. Výsledek \(3\cdot6!=2160\).
\end{prikladReseniBox}


\newpage
\blok{Samostatná práce a zásoba k procvičení}

U každé úlohy urči, zda se použijí všechny prvky a zda jsou místa či funkce rozlišeny. C5 ověřuje samostatné zopakování blokové metody.

% Zdroj C1: import_hodiny/32_Permutace_bez_opakovani.tex, Procvičování, položka 1.
\begin{zadaniReseniBox}[title={Úloha C1}]
Kolik pořadí má pět různých knih?
\tcblower
Řadíme všech pět různých knih, tedy $P(5)=5!=120$.
\end{zadaniReseniBox}

% Zdroj C2: import_hodiny/32_Permutace_bez_opakovani.tex, Procvičování, položka 2.
\begin{zadaniReseniBox}[title={Úloha C2}]
Kolika způsoby získají z deseti závodníků zlato a stříbro dva různí?
\tcblower
Zlato má deset možností, stříbro devět dalších. Role jsou rozlišené: $V(10,2)=10\cdot9=90$.
\end{zadaniReseniBox}

% Zdroj C3: import_hodiny/32_Permutace_bez_opakovani.tex, Procvičování, položka 3.
\begin{zadaniReseniBox}[title={Úloha C3}]
Kolik řad lze sestavit z pěti různých osob včetně Petry, je-li Petra vždy první?
\tcblower
Petřino místo je pevné, řadíme pouze čtyři ostatní osoby na čtyřech místech: $4!=24$.
\end{zadaniReseniBox}

% Zdroj C4: import_hodiny/32_Permutace_bez_opakovani.tex, Procvičování, položka 4.
\begin{zadaniReseniBox}[title={Úloha C4}]
Ze sedmi lidí obsadíme funkce vedoucího, zástupce a pokladníka třemi různými lidmi. Kolika způsoby?
\tcblower
Funkce nejsou zaměnitelné. Postupně volíme vedoucího sedmi způsoby, zástupce šesti a pokladníka pěti způsoby: $V(7,3)=7\cdot6\cdot5=210$.
\end{zadaniReseniBox}

% Zdroj C5: import_hodiny/32_Permutace_bez_opakovani.tex, Procvičování, položka 5.
\begin{zadaniReseniBox}[title={Úloha C5}]
Šest různých lidí postavíme do řady. Kolik je pořadí, musejí-li dva předem určení lidé stát vedle sebe?
\tcblower
Blok dvou lidí a čtyři zbývající osoby tvoří pět různých objektů. Uvnitř bloku jsou dvě pořadí, proto $5!\cdot2!=240$.
\end{zadaniReseniBox}


\poznamka{Typické chyby: použití $n!$, ačkoli obsazujeme jen dvě místa; zaměnění komise a funkcí; vynechání vnitřního pořadí bloku. U P5 začni omezenou poslední cifrou, teprve potom řaď ostatní. Příště odlišíme skutečně různé objekty od nerozlišitelných kopií a opravíme vícenásobné započtení.}
\end{document}
